\documentclass{article}

\usepackage{fullpage,amsmath}

\pagestyle{empty}

\begin{document}

\begin{center}
{\bf Calculus supplement: Periodic drug delivery}
\end{center}

In problem 10 in section 2.2 of Stewart, we consider a patient
receiving an injection of a drug at regular intervals.  Let's simplify
the problem to receiving a unit dose at intervals of 1 time unit.
One model for the delivery function is
\[
f(t) = \begin{cases}
e^{-t}, & 0 \leq t < 1 \\
e^{-(t-1)}, & 1 \leq t < 2 \\
e^{-(t-2)}, & 2 \leq t < 3 \\
\text{and so forth} \ldots 
\end{cases}
\]
where $e \approx 2.718\ldots$  Notice that it is
represented as a piecewise function where we ``glue'' together
segments of continuous curves.

\begin{enumerate}
\item Using a graphing
calculator as an aid, sketch this function.

\item One great mathematical achievement of the 19th century was the
  discovery that all periodic functions can be written as an \textit{infinite}
sum of
  sine and cosine functions.  These series are known as Fourier
  series.  The following are truncated (\textit{finite}) Fourier series for
$f(t)$.
\[
f_1(t) = \left(1-e^{-2}\right) \left(\frac{1}{2}
+ \frac{1}{1 + \pi^2}\cos(2 \pi t)  + 
\frac{\pi}{1 + \pi^2}\sin(2 \pi t)\right),
\]
\[
f_2(t) = \left(1-e^{-2}\right) \left(
\frac{1}{2} 
+ \frac{1}{1 + \pi^2}\cos(2 \pi t)  + 
\frac{\pi }{1 + \pi^2}\sin(2 \pi t) 
+ \frac{1}{1 + 4 \pi^2}\cos(4 \pi t)  + 
\frac{2 \pi}{1 + 4 \pi^2}\sin(4 \pi t)\right),
\]
\begin{multline*}
f_3(t) = \left(1-e^{-2}\right) \times \\
\left(
\frac{1}{2} 
+ \frac{1}{1 + \pi^2}\cos(2 \pi t)  + 
\frac{\pi }{1 + \pi^2}\sin(2 \pi t) 
+ \frac{1}{1 + 4 \pi^2}\cos(4 \pi t)  + 
\frac{2 \pi}{1 + 4 \pi^2}\sin(4 \pi t) \right. \\
\left. + \frac{1}{1 + 9 \pi^2}\cos(6 \pi t)  + 
\frac{3 \pi}{1 + 9 \pi^2}\sin(6 \pi t)\right),
\end{multline*}
\begin{multline*}
f_4(t) = \left(1-e^{-2}\right) \times \\
\left(
\frac{1}{2} 
+ \frac{1}{1 + \pi^2}\cos(2 \pi t)  + 
\frac{\pi }{1 + \pi^2}\sin(2 \pi t) 
+ \frac{1}{1 + 4 \pi^2}\cos(4 \pi t)  + 
\frac{2 \pi}{1 + 4 \pi^2}\sin(4 \pi t) \right. \\
\left. + \frac{1}{1 + 9 \pi^2}\cos(6 \pi t)  + 
\frac{3 \pi}{1 + 9 \pi^2}\sin(6 \pi t)
+ \frac{1}{1 + 16 \pi^2}\cos(8 \pi t)  + 
\frac{4 \pi}{1 + 16 \pi^2}\sin(8 \pi t)\right).
\end{multline*}
Using a graphing calculator as an aid, sketch $f(t)$, $f_1(t)$,
$f_2(t)$, $f_3(t)$, $f_4(t)$ on the same graph with the
domain $0 \leq t \leq 2$.  Are all the $f_n(t)$'s periodic?  What is
the period for each one?

\item Describe the relationship between $f(t)$ and the
$f_n(t)$'s.

\item Do you see a pattern in the Fourier series?  Write down an
  expression for $f_5(t)$.  

\item Determine the following quantities:
\begin{enumerate}

\item $\lim_{t \rightarrow 1^+} f(t)$.

\item $\lim_{t \rightarrow 1^-} f(t)$.

\item $\lim_{t \rightarrow 1^-} f_1(t)$.

\item $\lim_{t \rightarrow 1^+} f_1(t)$.

\item Generalize the left and right limits above to $f_n(t)$ for any $n$.

\end{enumerate}

\item Are the $f_n(t)$'s continuous?  Justify your answer using the limit
laws and the definition of continuity.  Is $f(t)$ continuous? Reconcile
  your answer with your answer to question 3.

\end{enumerate}
\end{document}
